\section{Lecture 11 (September 29th)}
\begin{ex}
(Effects of electron-electron repulsion) In the presence of the a potential, the electron-electron Coulomb repulsion can be treated as a pertubation. The energy shift is given as
\[\mathrm{\Delta }E=\iint d^3\vec{r}_{1}d^3\vec{r}_{2}\,\Big(\phi _{100}(\vec{r}_{1})\phi _{100}(\vec{r}_{2})\Big)^{*}\dfrac{e^2}{2\pi \varepsilon_0|\vec{r}_{1}-\vec{r}_{2}}\Big(\phi _{100}(\vec{r}_{1})\phi _{100}(\vec{r}_{2})\Big)\]
Notice that this reduces to
\[\mathrm{\Delta }E=\iint d^3\vec{r}_{1}d^3\vec{r}_{2}\,|\phi _{100}(\vec{r}_{1})|^2|\phi _{100}(\vec{r}_{2})|^2\dfrac{e^2}{2\pi \varepsilon _{0}|\vec{r}_{1}-\vec{r}_{2}|}\]
This can be worked out to give
\[\mathrm{\Delta }E=\dfrac{5}{8}\dfrac{Ze^2}{2\pi \varepsilon_0a_0}=\dfrac{5}{4}Z\Big(\dfrac{1}{2}mc^2\alpha ^2\Big)\]
as we know the zeroth order result to be $-108.8$ eV, we find up to first order, $E=-74.8$ eV.
\end{ex}
\vspace{2ex}
\begin{ex}
(First excited state) As we have two states with equal energies (spin triplet and spin singlet) we need to employ degenerate pertubation and diagonalise the $2\times 2$ matrix in the degenerate subspace. The matrix will look like, for
\[\psi _{\pm}=\dfrac{1}{\sqrt{2}}\Big(\phi _{100}(\vec{r}_{1})\phi _{2lm}(\vec{r}_{2})\pm \phi _{100}(\vec{r}_{2})\phi _{2lm}(\vec{r}_{1})\Big)\cdot X_{s,t}\]
The matrix is expected to look like
\[\begin{pmatrix}
\langle +|V_{12}|_{\rangle }&\langle +|V_{12}|-\rangle \\
\langle -|V_{12}|+\rangle &\langle -|V_{12}|-\rangle 
\end{pmatrix}
\]
However, the off-diagonal terms become zero and the matrix is already diagonal. We find the diagonal terms to be
\[\begin{align*}
\langle \pm |V_{12}|\pm\rangle =&\int d^3\vec{r}_{1}d^3\vec{r}_{2}\,\psi _{\pm}^{*}(\vec{r}_{1},\vec{r}_{2})\dfrac{e^2}{4\pi \varepsilon _{0}|\vec{r}_{1}-\vec{r}_{2}|}\psi _{\pm}(\vec{r}_{1},\vec{r}_{2})\\
=&J_{nl}+K_{nl}
\end{align*}\]
where $J_{nl}$ comes from the squares and $K_{nl}$ comes from the cross terms. 
\[\begin{cases}
\mathrm{\Delta }E^{(s)}_{nl}=J_{nl}+K_{nl}\\
\mathrm{\Delta }E^{(t)_{nl}}=J_{nl}-K_{nl}
\end{cases}\]
We find that the singlet state has a higher energy than the triplet state and that the spin had influence on the energy shift. We call this the spin exchange effect. Using Pauli matrices, we can write
\[\mathrm{\Delta }E=J_{nl}-\dfrac{1}{2}(1+\vec{\sigma }_{1}\cdot \vec{\sigma }_{2})K_{nl}\]
This develops to become Heisenberg's model of quantum magnetism. 
\end{ex}
\vspace{2ex}
\begin{rmk}
For the ground state, we discover that $(L=0,S=0,J=L+S=0)$. Using the spectroscopy notation ${}^{2S+1}L_{J}$ we find ${}^{1}S_{0}$ and we call this state parahelium. For the first excited state, we note that $(100)(2s)$ is lower in energy than the higher $(100)(2p)$ due to the orbital shape. Each are denoted ${}^{1}S_{0}$ and ${}^{3}S_{1}$ where the latter is called ortho helium.
\end{rmk}
\vspace{2ex}

